\begin{resumo}

\noindent % a NBR 6028:2021 dispensa o recuo de parágrafo para o resumo
A rede de distribuição de energia (PDN) de uma placa de circuito impresso
multicamadas é um dos principais determinantes da conformidade eletromagnética
do produto: ressonâncias de cavidade entre planos de alimentação e referência
elevam a impedância vista pelos circuitos integrados, convertendo ruído de
comutação em emissão radiada e comprometendo a aprovação em ensaios normativos.
A avaliação dessa impedância depende hoje de solvers eletromagnéticos de onda
completa, cujo custo computacional — dezenas de minutos por configuração —
inviabiliza a exploração sistemática do espaço de projeto nas fases iniciais do
layout, quando as decisões de baixo custo ainda são possíveis. Esta proposta
tem por objetivo desenvolver e validar um modelo substituto (\textit{surrogate})
informado por física capaz de predizer a curva completa de impedância
$Z(f)$ da PDN, no intervalo de 1\,MHz a 1\,GHz, a partir de oito parâmetros
geométricos e de material do empilhamento, com latência de inferência inferior
a um milissegundo. O trabalho utiliza o subconjunto de PDN de seis camadas da
SI/PI-Database da Universidade Técnica de Hamburgo, composto por 985
configurações amostradas por hipercubo latino e simuladas por método baseado em
física em 334 pontos de frequência. A contribuição central é a incorporação, à
função de perda da rede neural, de restrições eletromagnéticas disponíveis em
forma fechada — o comportamento capacitivo que rege o regime quase-estático, a
monotonicidade da resposta até o nulo de série e a condição de passividade do
sistema — de modo a compensar a escassez de amostras, regime em que modelos
puramente orientados a dados tendem a violar a física do problema. Uma
investigação preliminar já conduzida sobre a base confirmou a validade da âncora
capacitiva, cuja inclinação medida vale $-1{,}018 \pm 0{,}008$ contra $-1$
teórico, e afastou a ancoragem em frequências modais, não observáveis na
autoimpedância dentro da faixa analisada. O modelo será comparado a uma bateria de
referências (regressão linear regularizada, florestas aleatórias, gradient
boosting, processos gaussianos e perceptron multicamadas), integrado a um
otimizador evolutivo multiobjetivo para busca de empilhamentos de menor risco
de interferência eletromagnética, e interpretado por regressão simbólica e
valores de Shapley, com o objetivo de destilar diretrizes de projeto legíveis
por engenheiros. Espera-se entregar, ao final, um modelo com coeficiente de
determinação superior a 0,90 na predição da impedância máxima em conjunto de
teste, um conjunto de expressões analíticas interpretáveis relacionando
parâmetros de empilhamento e risco eletromagnético, e uma implementação
reprodutível de código aberto documentada segundo práticas de desenvolvimento
orientado a especificação.

\end{resumo}
